\section{Introducción}
El presente informe de implementación documenta de manera exhaustiva el proceso de adopción, configuración y puesta en marcha del ecosistema de gestión en la nube basado en la plataforma Odoo para la empresa Perú Flores Travel Agency Tours Operador E.I.R.L.. El objetivo central de este proyecto es ejecutar una transformación digital integral que permita optimizar los flujos comerciales y operativos, centralizar la información y erradicar las ineficiencias derivadas de un modelo de gestión tradicional.

Perú Flores es una agencia de viajes con amplia experiencia en el turismo receptivo, especializada en organizar y operar paquetes turísticos personalizados hacia destinos clave como Cusco, Arequipa y Puno. Con un fuerte enfoque comercial hacia el mercado extranjero, especialmente el turista brasileño, la agencia se ha sostenido gracias a la recomendación directa y al esfuerzo de su gerencia general. Sin embargo, toda esta operación se ha gestionado de manera 100\% manual. Actualmente, la captación de leads y la comunicación se centralizan informalmente en WhatsApp, mientras que el registro de clientes se realiza en cuadernos físicos y la elaboración de cotizaciones depende de copiar, pegar y modificar manualmente antiguos documentos de Word y plantillas de Excel creadas durante la pandemia.

A pesar de la calidad de sus servicios, este modelo fragmentado ha generado cuellos de botella críticos. La gerencia invierte una cantidad excesiva de tiempo realizando cálculos manuales repetitivos (como dividir el costo de guías y transportes entre el número de pasajeros para armar un paquete), y calculando márgenes de ganancia e impuestos (IGV) de forma artesanal para cada prospecto. A esto se suma la duplicidad de costos tecnológicos, al mantener pagos anuales por servicios de hosting web que no están integrados con la operación diaria, y la incapacidad de recopilar testimonios automatizados que potencien el ``boca a boca'' digital tras finalizar un servicio.

En este contexto, la implementación de Odoo se plantea no solo como un cambio de software, sino como una solución estratégica para profesionalizar la atención al turista y ordenar el flujo logístico y comercial. Mediante el despliegue de los módulos de CRM, Ventas, Sitio Web y Encuestas, la plataforma permitirá a Perú Flores obtener una base de datos centralizada, generar cotizaciones en formato PDF de manera automatizada, y medir la satisfacción del viajero mediante formularios digitales estructurados. Todo esto, eliminando la necesidad de mantener infraestructuras web externas y costosas.

En coherencia con la planificación general y las lecciones aprendidas de consultorías previas, este proyecto se rige bajo un enfoque estrictamente Out-of-the-box (uso de funcionalidades nativas). Se ha descartado cualquier desarrollo de código a medida para asegurar una adopción ágil, rápida y libre de riesgos técnicos. Asimismo, se establecen límites claros: la facturación electrónica seguirá gestionándose externamente a través del portal de la SUNAT y no se realizarán integraciones complejas mediante APIs con terceros.

En síntesis, esta introducción sitúa el punto de partida del proyecto: la transición de Perú Flores desde una gestión basada en la memoria, los cuadernos y las sumas manuales, hacia un ecosistema unificado en la nube que le otorgue agilidad comercial, reduzca su carga administrativa y le permita escalar su nivel de servicio en el competitivo mercado turístico.
